\chapter{Ciąg dokładny chirurgii i zbiór struktur}
\label{ch:ciag-chirurgii}
\index{ciag dokladny chirurgii@Ciąg dokładny chirurgii}

\noindent\textbf{Pytanie przewodnie.}
Twierdzenie~\ref{thm:Wall-L-22} rozstrzyga, czy
\emph{jedno} normalne odwzorowanie można poprawić
do prostej równoważności homotopijnej.
Do klasyfikacji potrzebujemy jeszcze wiedzieć,
które odwzorowania normalne liczymy i kiedy
dwie otrzymane równoważności opisują tę samą
gładką rozmaitość. Odpowiedź łączy normalne
niezmienniki, grupy $L$ i zbiór struktur.
W całym rozdziale, o ile nie zaznaczono inaczej,
$X^n$ jest spójną, zamkniętą, gładką
rozmaitością wymiaru $n\geq5$.
Ta konwencja pozwala stosować gładką
wersję Twierdzenia~\ref{thm:s-kobordyzm-20}
do kobordyzmów wymiaru $n+1\geq6$.

\section{Od normalnego odwzorowania do normalnego niezmiennika}
\label{sec:normalne-niezmienniki-23}

Definicja~\ref{def:normal-map-21} wiąże
odwzorowanie $f:M\to X$ z izomorfizmem
stabilnych wiązek normalnych nad pewną
wiązką $\xi\to X$. Dwie takie trójki
uważamy za tę samą \emph{daną normalną},
gdy łączy je normalny kobordyzm z
Definicji~\ref{def:normal-bordism-21}.
Nie żądamy jeszcze, aby $f$ było
równoważnością homotopijną.

\begin{definicja}[Zbiór normalnych niezmienników]
\index{zbior normalnych niezmiennikow@Zbiór normalnych niezmienników}
\label{def:normal-invariants-23}
Przez $\mathcal N_{\mathrm{DIFF}}(X)$ oznaczamy
zbiór klas normalnych odwzorowań stopnia jeden
$(f,b,\xi):M^n\to X^n$, gdzie wiązka
$\xi\to X$ może się zmieniać, względem normalnego
kobordyzmu. Punktem wyróżnionym jest
$(\id_X,\id_{\nu_X},\nu_X)$.
\end{definicja}

\begin{przyklad}[Identyczność]
\label{prz:normalny-punkt-23}
Identyczność $X\to X$ ma stopień jeden i
kanoniczny izomorfizm wiązek normalnych
przy wyborze $\xi=\nu_X$.
Stanowi zero zbioru normalnych niezmienników.
Z kolei odwzorowanie stopnia jeden bez
izomorfizmu~\eqref{eq:normal-map-21}
nie określa elementu $\mathcal N_{\mathrm{DIFF}}(X)$;
sam stopień nie pamięta obramowań
potrzebnych do chirurgii.
\end{przyklad}

Aby zapisać te dane homotopijnie, niech
$O=\operatorname*{colim}_{q}O(q)$, a $G$ oznacza
stabilną przestrzeń samorównoważności
homotopijnych sfer (składowe stopnia $\pm1$).
Odwzorowanie $O\to G$ zapomina liniową
strukturę obrotu sfery. Jego stabilny iloraz
homotopijny oznaczamy $G/O$.
W równoważnym opisie jest on włóknem
homotopijnym $BO\to BG$: mierzy wybór
stabilnej struktury wektorowej na zadanej
stabilnej sferycznej fibracji.
Na przykład dwie konstrukcje z tą samą
fibracją sferyczną mogą mieć różne
stabilne wiązki wektorowe; różnicę rejestruje
właśnie klasa w $G/O$.

\begin{twierdzenie}[Interpretacja normalnych niezmienników;
 wynik zewnętrzny]
\label{thm:normal-invariants-GO-23}
Dla zamkniętej gładkiej rozmaitości $X$
konstrukcja Pontriagina--Thoma daje
naturalną bijekcję zbiorów z punktami
wyróżnionymi
\[
 \mathcal N_{\mathrm{DIFF}}(X)\cong [X,G/O].
\]
Prawa strona oznacza klasy homotopii
odwzorowań do $G/O$; punkt stały odpowiada
identyczności $X\to X$.
\end{twierdzenie}

\noindent\textbf{Co stoi za bijekcją.}
Zanurzamy $M$ w wysokowymiarowej przestrzeni
euklidesowej i zwijamy dopełnienie jej
otoczenia rurkowego do punktu.
Izomorfizm $b$ wraz z warunkiem stopnia jeden
porównuje otrzymaną fibrację sferyczną z tą
nad $X$; wybór $\xi$ może się zmieniać.
Kobordyzm normalny staje się homotopią
wyborów podniesienia do $BO$.
W drugą stronę transwersalna reprezentacja
klasy $[X,G/O]$ daje podrozmaitość i
odwzorowanie stopnia jeden.
Twierdzenie transwersalności i stabilny
rachunek Pontriagina--Thoma są tutaj
wynikami zewnętrznymi; pełny dowód
identyfikacji znajduje się u
\href{https://webhomes.maths.ed.ac.uk/~v1ranick/books/scm.pdf}
{Walla, \emph{Surgery on Compact Manifolds},
stwierdzenia 10.1–10.2 i lemat 10.6}.
Stwierdzenie 10.2 utożsamia klasy normalnych
map z danymi redukcji, a lemat 10.6
identyfikuje te dane z $[X,G/O]$.
Ten opis wyjaśnia intuicję, nie zastępuje
dowodu konstrukcji odwrotnej.

\noindent\textbf{Jak czytać klasę w $G/O$.}
Niech $\nu_X^{\mathrm{sph}}:X\to BG$ klasyfikuje
stabilne rozwłóknienie sferyczne wyznaczone
przez wiązkę normalną $X$.
Dana normalna wybiera nad nim stabilną
wiązkę wektorową $\xi$. Jej odwzorowanie
klasyfikujące $c_\xi:X\to BO$ daje
podniesienie w poniższym diagramie,
wraz z homotopią
jego złożenia do $BG$ z $\nu_X^{\mathrm{sph}}$:
\[
\begin{tikzcd}[column sep=large]
 & BO \arrow[d] \\
 X \arrow[ur,dashed,"c_\xi"]
   \arrow[r,"\nu_X^{\mathrm{sph}}"'] & BG .
\end{tikzcd}
\]
Wiązka $\nu_X$ daje podniesienie wyróżnione.
Porównanie z nim wyznacza klasę w $[X,G/O]$.
Przejście od trójki $(f,b,\xi)$ do tego
podniesienia jest częścią przywołanej
identyfikacji Pontriagina--Thoma;
sam diagram nie konstruuje $f$.

\section{Jak rozróżnić gładkie realizacje celu?}
\label{sec:zbior-struktur-23}

Wcześniej, w Rozdziale~\ref{ch:s-kobordyzm},
odróżniliśmy zwykłą od prostej równoważności.
To rozróżnienie zachowujemy w klasyfikacji.

\begin{definicja}[Prosty gładki zbiór struktur]
\index{prosty gladki zbior struktur@Prosty gładki zbiór struktur}
\label{def:structure-set-23}
$\mathcal S^s_{\mathrm{DIFF}}(X)$ jest zbiorem
klas par $(M,f)$, gdzie $M$ jest zamkniętą
gładką rozmaitością, a $f:M\to X$ prostą
równoważnością homotopijną. Pary $(M_0,f_0)$
i $(M_1,f_1)$ są równoważne, gdy istnieje
dyfeomorfizm $h:M_0\to M_1$ taki, że
$f_1\circ h\simeq f_0$. Punktem wyróżnionym
jest $(X,\id_X)$.
\end{definicja}

Intuicja: element $\mathcal S^s_{\mathrm{DIFF}}(X)$
to gładka realizacja ustalonego typu homotopijnego
razem z identyfikacją tego typu z $X$.

\begin{przyklad}[Ten sam typ homotopijny]
\label{prz:struktury-sfera-23}
Jeśli $\Sigma^n$ jest gładką sferą homotopijną,
każda wybrana prosta równoważność
$\Sigma\to S^n$ daje element
$\mathcal S^s_{\mathrm{DIFF}}(S^n)$.
Samo stwierdzenie $\Sigma\simeq S^n$
nie orzeka, czy ten element jest punktem
wyróżnionym: definicja wymaga dyfeomorfizmu.
Konkretne niezerowe przykłady zobaczymy
w Rozdziale~\ref{ch:sfery-Milnora}.
\end{przyklad}

Prosta równoważność $f$ ma kanoniczną klasę
normalną: przenosi stabilną normalną
fibrację sferyczną, a stabilną wiązkę
$\nu_M$ można przenieść na wiązkę $\xi$
nad $X$ za pomocą odwrotności homotopijnej $f$.
Tak otrzymana wiązka nie musi być
stabilnie izomorficzna z $\nu_X$.
Wybór wyznacza element $\mathcal N_{\mathrm{DIFF}}(X)$.
Otrzymujemy odwzorowanie
\[
 \eta:\mathcal S^s_{\mathrm{DIFF}}(X)
          \longrightarrow\mathcal N_{\mathrm{DIFF}}(X).
\]
To nie jest dowolne zapomnienie:
jeżeli dwie struktury mają to samo $\eta$,
istnieje między nimi normalny kobordyzm
nad $X\times I$. Pozostała przeszkoda
tego kobordyzmu leży w grupie $L_{n+1}$.
Konkretnie $\eta([M,f])$ pamięta stabilną
wiązkę normalną $M$ przeniesioną na $X$
przez odwrotność homotopijną $f$ i jej
porównanie ze sferyczną wiązką normalną $X$.
Nie pamięta natomiast samej prostej
równoważności $f$ jako struktury gładkiej.

\section{Działanie grupy \texorpdfstring{$L$}{L} i ciąg dokładny}
\label{sec:exact-surgery-23}

Niech $\pi=\pi_1(X)$, a
$w:\pi\to\{\pm1\}$ będzie charakterem
orientacji. Twierdzenie~\ref{thm:Wall-L-22}
określa odwzorowanie
\[
 \theta:\mathcal N_{\mathrm{DIFF}}(X)
       \longrightarrow L_n^s(\Z[\pi],w).
\]
Zbiór $[X,G/O]$ ma naturalną strukturę grupy
abelowej, lecz $\theta$ na ogół nie jest
homomorfizmem tej grupy. W ciągu używamy go
jako odwzorowania zbiorów z punktami wyróżnionymi;
do pojęcia dokładności wystarcza przeciwobraz zera.

\begin{twierdzenie}[Ciąg dokładny chirurgii;
 wynik Walla]
\label{thm:exact-surgery-23}
Dla $X$ jak na początku rozdziału istnieje
ciąg zbiorów z punktami wyróżnionymi
\begin{equation}\label{eq:exact-surgery-23}
 [\Sigma X_+,G/O]\xrightarrow{\ \theta_{n+1}\ }
 L_{n+1}^s(\Z[\pi],w)
 \xrightarrow{\ \partial_{\!s}\ }
 \mathcal S^s_{\mathrm{DIFF}}(X)
 \xrightarrow{\ \eta\ }
 [X,G/O]
 \xrightarrow{\ \theta_n\ }
 L_n^s(\Z[\pi],w).
\end{equation}
Tutaj $X_+$ oznacza $X$ z dodanym punktem
bazowym, a $\Sigma$ zawieszenie zredukowane.
W odróżnieniu od $\theta_n$ lewa strzałka
$\theta_{n+1}$ jest homomorfizmem grup;
to osobny wniosek z konstrukcji na cylindrze
(Wall, stwierdzenie~10.7).
Grupa $L_{n+1}^s$ działa na zbiorze struktur;
$\partial_{\!s}$ wysyła jej element na wynik
działania na strukturze $(X,\id_X)$.
Dokładność oznacza:
\begin{enumerate}[label=(\roman*)]
\item obraz $\eta$ jest dokładnie
      $\theta_n^{-1}(0)$;
\item dwie struktury mają ten sam obraz
      przez $\eta$ wtedy i tylko wtedy,
      gdy leżą w jednej orbicie działania
      $L_{n+1}^s$;
\item stabilizator punktu wyróżnionego
      w $L_{n+1}^s$ jest obrazem
      $\theta_{n+1}$.
\end{enumerate}
\end{twierdzenie}

\begin{proof}[Wyjaśnienie kolejnych miejsc dokładności]
Najpierw weźmy normalną klasę $a$
z $\theta_n(a)=0$. Twierdzenie Walla
daje normalny kobordyzm od jej
reprezentanta do prostej równoważności.
Ta równoważność określa strukturę
z $\eta([M,f])=a$. Odwrotnie struktura
jest już prostą równoważnością,
więc jej przeszkoda znika.
To dowodzi (i).

Jeśli $\eta([M_0,f_0])=\eta([M_1,f_1])$,
normalny kobordyzm między dwiema
strukturami jest problemem chirurgicznym
nad $X\times I$. Jego przeszkoda ma
wymiar $n+1$ i należy do $L_{n+1}^s$.
Geometryczne twierdzenia realizacyjne
Walla pozwalają zrealizować każdy taki
element przez odpowiedni kobordyzm:
dla parzystego $n+1\geq6$ jest to
twierdzenie 5.8, a dla nieparzystego
$n+1\geq7$ twierdzenie 6.5.
W szczególności przy $n=5$ używamy
parzystego twierdzenia 5.8 dla $L_6^s$;
twierdzenie 6.5 zaczyna się od wymiaru siedem.
Operacja na górnym brzegu daje działanie.
Aby porównać jej wynik z $[M_1,f_1]$,
oznaczmy dany ślad $M_0\to M_1$ przez $W$,
a ślad realizujący tę samą przeszkodę
$\alpha$ od $M_0$ do $M_\alpha$ przez $W_\alpha$.
Sklejamy odwrócony $W_\alpha$ z $W$
wzdłuż $M_0$. Addytywność daje
przeszkodę $-\alpha+\alpha=0$ dla śladu
$M_\alpha\to M_1$. Chirurgia relatywnie do końców
zamienia je w prostą równoważność nad
$X\times I$. Ponieważ oba końce są już
prostymi równoważnościami z $X$,
ich inkluzje do otrzymanego kobordyzmu
są prostymi równoważnościami; powstaje
$s$-kobordyzm, a
Twierdzenie~\ref{thm:s-kobordyzm-20}
zamienia go w dyfeomorfizm końców.
W drugą stronę działanie zachowuje
normalną klasę, bo jego ślad jest
normalnym kobordyzmem. To daje (ii).

Jeśli element $L_{n+1}^s$ pochodzi
z normalnej klasy na cylindrze
$X\times I$ ustalonej na obu końcach,
jego działanie na $(X,\id_X)$
zwraca tę samą strukturę.
Twierdzenie o realizacji i
odwrotne rozpoznanie takiego
cylindra dają również przeciwną
implikację; jest to (iii).
Utożsamienie normalnych danych cylindra
z $[\Sigma X_+,G/O]$ wynika
z ilorazu $X\times I/(X\times\partial I)$.
Te głębokie kroki są odpowiednio
twierdzeniami 10.4, 10.5 i 10.8 w
\href{https://webhomes.maths.ed.ac.uk/~v1ranick/books/scm.pdf}
{książce Walla, §10}; przywołaliśmy
ich hipotezy i zastosowanie, a nie
powtórzyliśmy ich pełnego dowodu.
\end{proof}

\noindent\textbf{Co robią cztery strzałki.}
Każda z nich ma inne dane wejściowe;
strzałki przy zbiorze struktur nie są
homomorfizmami grup.
\begin{enumerate}[label=(\roman*)]
\item $\theta_{n+1}$ bierze normalną daną
      na $X\times I$ ustaloną na obu końcach.
      Po ich sklejeniu w punkt jest to klasa
      $[\Sigma X_+,G/O]$. Przeszkoda chirurgiczna
      tego $(n+1)$-wymiarowego problemu jest
      elementem $L_{n+1}^s$. Sklejanie cylindrów
      odpowiada dodawaniu przeszkód, co wyjaśnia
      homomorficzność tej strzałki.
\item $\partial_{\!s}$ zaczyna od
      $\alpha\in L_{n+1}^s$.
      Twierdzenie realizacyjne Walla dostarcza
      normalnego śladu $W\to X\times I$:
      na dole mamy $(X,\id_X)$, a na górze
      prostą równoważność $(M_\alpha,f_\alpha)$.
      Wynik to $[M_\alpha,f_\alpha]$.
      Tę konstrukcję, wykonaną z dowolną
      strukturą na dole, stosujemy jako działanie
      $\alpha$ na całym
      $\mathcal S^s_{\mathrm{DIFF}}(X)$.
\item $\eta$ wysyła $[M,f]$ na klasę
      normalnego odwzorowania indukowaną
      przez $f$; jej równość dla dwóch struktur
      oznacza istnienie normalnego kobordyzmu
      między reprezentantami.
\item $\theta_n$ przypisuje normalnej klasie
      przeszkodę jej przedstawiciela.
      Po chirurgii poniżej środka wymiaru
      mierzy ją kwadratowe jądro w wymiarze
      parzystym albo formacja w nieparzystym,
      zgodnie z Rozdziałem~\ref{ch:przeszkoda-chirurgiczna}.
      Jej znikanie jest warunkiem poprawienia
      mapy normalnej do prostej równoważności.
\end{enumerate}
Rozkład na te operacje wyjaśnia też dwa
różne wymiary grup $L$: obiekt nad $X$
ma wymiar $n$, a każdy ślad nad $X\times I$
ma wymiar $n+1$.

\begin{figure}[htbp]
\centering
\begin{tikzpicture}[>=Latex,font=\small,
 every node/.style={align=center}]
 \node[draw=manifoldblue!75!black,fill=manifoldblue!8,
   rounded corners,minimum width=1.9cm,minimum height=.7cm]
   (m0) at (0,2) {$M_0$};
 \node[draw=manifoldblue!75!black,fill=manifoldblue!8,
   rounded corners,minimum width=2.5cm,minimum height=.7cm]
   (w) at (3.6,2) {$W^{n+1}$};
 \node[draw=manifoldblue!75!black,fill=manifoldblue!8,
   rounded corners,minimum width=1.9cm,minimum height=.7cm]
   (m1) at (7.2,2) {$M_1$};
 \node (x0) at (0,0) {$X\times\{0\}$};
 \node (xi) at (3.6,0) {$X\times I$};
 \node (x1) at (7.2,0) {$X\times\{1\}$};
 \draw[->,thick,manifoldblue!70!black]
   (m0)--node[above] {$i_0$} (w);
 \draw[<-,thick,manifoldblue!70!black]
   (w)--node[above] {$i_1$} (m1);
 \draw[->,thick,manifoldblue!70!black]
   (m0)--node[left] {$f_0$} (x0);
 \draw[->,thick,manifoldblue!70!black]
   (w)--node[right] {$F$} (xi);
 \draw[->,thick,manifoldblue!70!black]
   (m1)--node[right] {$f_1$} (x1);
 \draw[->,thick,manifoldblue!70!black]
   (x0)--(xi);
 \draw[<-,thick,manifoldblue!70!black]
   (xi)--(x1);
\end{tikzpicture}
\caption{Normalny ślad nad cylindrem $X\times I$.
Dolny brzeg $M_0$ jest daną początkową,
górny $M_1$ wynikiem działania. Przy lewej
strzałce $\theta_{n+1}$ oba końce są ustalone
jako $X$; przy działaniu $L_{n+1}^s$
górny koniec może się zmienić.}
\label{fig:surgery-cobordism-cylinder-23}
\end{figure}

W rysunku strzałki poziome są inkluzjami
brzegów, a kwadraty komutują z dokładnością
do homotopii zgodnej z danymi normalnymi.
Przy $\eta(s_0)=\eta(s_1)$ istnieje taki
normalny ślad między strukturami.
Jego przeszkoda określa element przenoszący
$s_0$ w $s_1$, więc włókno $\eta$ jest orbitą.
Gdy oba końce reprezentują punkt wyróżniony,
możemy wybrać dyfeomorfizm identyfikujący
górną strukturę z $(X,\id_X)$.
Utożsamienie takiego śladu z normalną daną
cylindra i odwrotna implikacja są treścią
Twierdzenia 10.8 Walla: obraz $\theta_{n+1}$
jest dokładnie stabilizatorem bazy.

\clearpage
\section{Wersja z ustalonym brzegiem}
\label{sec:relative-surgery-23}
\index{ciag dokladny chirurgii@Ciąg dokładny chirurgii!wersja z ustalonym brzegiem@wersja z ustalonym brzegiem}

Niech teraz $X^n$ będzie zwartą, spójną,
gładką rozmaitością wymiaru $n\geq5$
z niepustym brzegiem. Ustalamy jego
gładką strukturę i identyfikację brzegu
każdej rozważanej rozmaitości z $\partial X$.
Przez $\mathcal S^s_{\mathrm{DIFF}}(X\operatorname{rel}\partial X)$
rozumiemy klasy prostych równoważności par
$f:(M,\partial M)\to(X,\partial X)$,
których ograniczenie do brzegu jest tą
ustaloną dyfeomorficzną identyfikacją.
Równoważność par wymaga dyfeomorfizmu
zgodnego z nią na brzegu i homotopii
odwzorowań ustalonej na brzegu.
Normalne odwzorowania i ich kobordyzmy
mają na kołnierzu brzegu dane produktowe.
Stopień jeden rozumiemy względem
względnych klas fundamentalnych:
$f^*w_X=w_M$ oraz
$f_*[M,\partial M]=[X,\partial X]$
w $H_n(X,\partial X;\Z^{w_X})$.
Intuicja jest taka sama jak poprzednio,
ale żadna chirurgia nie zmienia już
zadanej gładkiej struktury na $\partial X$.

\begin{twierdzenie}[Ciąg chirurgii z ustalonym brzegiem;
 wynik zewnętrzny]
\label{thm:relative-surgery-23}
Niech $\pi=\pi_1(X)$, a $w$ będzie jego
charakterem orientacji. Przy powyższych
hipotezach i konwencji o ustalonym brzegu,
jeśli $Q=X/\partial X$, istnieje ciąg
zbiorów z punktami wyróżnionymi
\begin{equation}\label{eq:relative-surgery-23}
 [\Sigma Q,G/O]\xrightarrow{\ \theta_{n+1}^{\mathrm{rel}}\ }
 L_{n+1}^s(\Z[\pi],w)
 \xrightarrow{\ \partial_{\!s}^{\mathrm{rel}}\ }
 \mathcal S^s_{\mathrm{DIFF}}(X\operatorname{rel}\partial X)
 \xrightarrow{\ \eta^{\mathrm{rel}}\ }
 [Q,G/O]
 \xrightarrow{\ \theta_n^{\mathrm{rel}}\ }
 L_n^s(\Z[\pi],w).
\end{equation}
Dokładność ma to samo znaczenie orbit,
stabilizatora i przeciwobrazu zera co
w Twierdzeniu~\ref{thm:exact-surgery-23}.
Nie zakładamy, by inkluzje składowych
brzegu indukowały izomorfizmy grup
podstawowych.
\end{twierdzenie}

Iloraz $Q$ zapamiętuje warunek na brzegu:
odwzorowania $X\to G/O$ wysyłające
$\partial X$ w punkt bazowy, z homotopiami
ustalonymi na $\partial X$, odpowiadają
opartym odwzorowaniom $Q\to G/O$.
Analogicznie normalne dane na $X\times I$
ustalone na $X\times\partial I$ i
$\partial X\times I$ są klasyfikowane przez
$[\Sigma Q,G/O]$. Ściana $\partial X\times I$
na rysunku~\ref{fig:surgery-cobordism-cylinder-23}
pozostaje wtedy cylindrem produktowym.
Wyniki z
\href{https://webhomes.maths.ed.ac.uk/~v1ranick/books/scm.pdf}
{§10 książki Walla, \emph{Surgery on Compact Manifolds}},
w tym gładki odpowiednik Twierdzenia 10.8,
dają identyfikację danych normalnych,
kryterium realizacji, działanie i dokładność
również w tej wersji. To są użyte tu
wyniki zewnętrzne, a nie wniosek wyłącznie
z zastąpienia $X$ przez $Q$ w poprzednim wzorze.

\begin{przyklad}[Dysk z ustalonym brzegiem]
Dla $X=D^n$ mamy $Q=D^n/S^{n-1}\cong S^n$.
Względne dane normalne są więc klasami
$[S^n,G/O]=\pi_n(G/O)$, a dane cylindra
klasami $\pi_{n+1}(G/O)$.
To pokazuje, dlaczego warunek na brzegu
jest istotny: bez niego $[D^n,G/O]$
ma tylko klasę stałą.
Przykład identyfikuje dane wejściowe ciągu;
nie oblicza jeszcze jego zbioru struktur.
\end{przyklad}

\section{Co ciąg mówi w praktyce?}
\label{sec:applications-surgery-23}

\begin{przyklad}[Cel po prostu spójny]
\label{prz:simply-connected-surgery-23}
Gdy $\pi=1$, obliczenie z
Sekcji~\ref{sec:obliczenia-LZ-22} daje
\[
 L_{4j}^s(\Z)\cong\Z,\qquad
 L_{4j+2}^s(\Z)\cong\Z/2,\qquad
 L_{2j+1}^s(\Z)=0.
\]
W wymiarze $4j$ pierwsza przeszkoda
normalnego odwzorowania jest sygnaturą
formy jądra podzieloną przez $8$.
Rozszczepienie z
Lematu~\ref{lem:split-normalne-21}
jest ortogonalne dla przecięć. Oznaczmy przez
$K=\ker\bigl(f_*:H_{2j}(M;\Q)\to
H_{2j}(X;\Q)\bigr)$. Dla $y\in K$
i $x\in H_{2j}(X;\Q)$
naturalność przecięć daje
$\lambda_M(f_!x,y)=\lambda_X(x,f_*y)=0$.
Dlatego sygnatura jądra to
$\operatorname{sign}(M)-\operatorname{sign}(X)$.
W wymiarze $4j+2$ mierzy ją
niezmiennik Arfa formy kwadratowej.
W wymiarach nieparzystych grupa przeszkód
jest zerowa. W każdym z tych przypadków
ciąg nadal wymaga obliczenia $[X,G/O]$
oraz działania $L_{n+1}$, zanim
rozróżnimy wszystkie struktury.
\end{przyklad}

\clearpage
\begin{przyklad}[Ślad sfery i produktu w ciągu chirurgii]
\label{prz:slad-s6-ciag-23}
Niech $W$ i $(f,b):M=S^3\times S^3\to S^6$
będą śladem i jego górną mapą normalną z
Przykładu~\ref{prz:slad-s6-produkt-21}.
Od strony $S^6$ ślad ma uchwyt indeksu trzy;
od strony $M$ uchwyt indeksu cztery na
$B=\{0\}\times S^3$ w nowo dołączonej
połowie $D^3\times S^3$ przywraca $S^6$.
To ta sama operacja oglądana z dwóch końców.

Połóżmy $A=S^3\times\{*\}$. Ponieważ
$H_3(S^6)=0$, jądro górnej mapy jest
$K_3(f)=H_3(M)=\Z[A]\oplus\Z[B]$.
Przykład~\ref{prz:forma-produkt-21} oblicza
$\lambda(A,B)=1$ oraz
$\lambda(A,A)=\lambda(B,B)=0$.
Lemat~\ref{lem:wall-lokalna-stabilizacja-22}
oraz Przykład~\ref{prz:produkt-L-22} dają
$\mu(A)=\mu(B)=0$ dla standardowego wyboru
$b$ w tym śladzie. Zatem dla
$\bar x,\bar y\in\mathbb F_2$ funkcja
kwadratowa ma postać
\[
 q(\bar x[A]+\bar y[B])=\bar x\bar y,
 \qquad \operatorname{Arf}(q)=0.
\]
Według rachunku z
Sekcji~\ref{sec:obliczenia-LZ-22}
przeszkoda tej mapy jest więc zerem
w $L_6^s(\Z)\cong\Z/2$.

Normalny ślad $W$ łączy $(f,b)$ z
$(S^6,\id_{S^6})$. Dlatego klasa normalna
$a=[f,b]\in[S^6,G/O]$ jest punktem wyróżnionym:
$a=0$. W ciągu~\eqref{eq:exact-surgery-23}
otrzymujemy konkretnie
\[
 [S^6,\id_{S^6}]
 \xrightarrow{\ \eta\ }a=0
 \xrightarrow{\ \theta_6\ }0\in L_6^s(\Z).
\]
Ponieważ $L_7^s(\Z)=0$ z~\eqref{eq:L-Z-22},
opis orbit w
Twierdzeniu~\ref{thm:exact-surgery-23} mówi,
że standardowa sfera jest jedyną strukturą
w tym włóknie $\eta$.
Sama mapa $f$ nie jest elementem zbioru
struktur: jej jądro $K_3(f)$ ma rangę dwa,
więc $f$ nie jest równoważnością homotopijną.
Wniosek $a=0$ korzysta z konkretnych danych
normalnych wybranych przez $W$.
\end{przyklad}

\begin{stwierdzenie}[Sfery homotopijne jako struktury nad $S^7$]
\label{prop:structure-set-seven-sphere-23}
\index{zbior struktur@Zbiór struktur!sfera wymiaru siedem@sfera wymiaru siedem}
Istnieje bijekcja zbiorów z punktami wyróżnionymi
\[
 \mathcal S^s_{\mathrm{DIFF}}(S^7)\cong\Theta_7,
\]
gdzie $\Theta_7$ oznacza zorientowane klasy
dyfeomorfizmu gładkich sfer homotopijnych.
\end{stwierdzenie}
\begin{proof}
Ustalmy orientację $S^7$.
Para $(M,f)$ ze zbioru struktur nadaje $M$
orientację, dla której $f:M\to S^7$ ma stopień
$+1$. Ponieważ $f$ jest równoważnością
homotopijną, $M$ jest sferą homotopijną.
Odwrotnie każda zorientowana sfera homotopijna
ma równoważność homotopijną stopnia $+1$
z $S^7$. Jest ona prosta, bo $\pi_1(S^7)=1$
i $\operatorname{Wh}(1)=0$.
Dwie mapy stopnia $+1$ z tej samej
siedmiowymiarowej sfery homotopijnej
do $S^7$ są homotopijne: ponieważ $S^7$
jest $6$-spójna, klasy takich map z
siedmiowymiarowego kompleksu CW
wyznacza stopień. Wobec tego zgodność
orientacji przy dyfeomorfizmie dokładnie
odpowiada równoważności par ze zbioru
struktur. Klasa $S^7$ odpowiada punktowi
wyróżnionemu.
\end{proof}

Twierdzenie~\ref{thm:all-seven-spheres-25}
podaje $28$ elementów $\Theta_7$, w tym
klasy sum spójnych sfery Milnora $M_3$.
To daje konkretną interpretację zbioru
struktur w ciągu~\eqref{eq:exact-surgery-23}.
Liczba $28$ wynika tam z wyników
Kervaire'a--Milnora i Eellsa--Kuipera;
niniejszy rozdział nie oblicza obrazu
lewej strzałki $\theta_8$ i nie wyprowadza
tej liczby z samego ciągu.

\begin{przyklad}[Dlaczego zerowa przeszkoda
nie oznacza jedynej struktury]
\label{prz:nonunique-structure-23}
Dla $\theta_n(a)=0$ istnieje
\emph{co najmniej jedna} prosta równoważność
o normalnej klasie $a$. Wszystkie
struktury z tą klasą tworzą jedną orbitę
$L_{n+1}^s$, która może zawierać
więcej niż jeden element.
To jest algebraiczne miejsce na różne
gładkie realizacje tego samego typu
homotopijnego. Sfery Milnora z
Rozdziału~\ref{ch:sfery-Milnora}
dadzą konkretny powód, dla którego
samej równoważności homotopijnej
nie wolno utożsamiać z dyfeomorfizmem.
\end{przyklad}

\noindent\textbf{Granica tej części.}
Ciąg~\eqref{eq:exact-surgery-23} kończy
program \emph{wysokowymiarowej} chirurgii
w sensie logicznym: podaje kryterium
istnienia struktury i opis niejednoznaczności.
Nie jest automatyczną tabelą dyfeomorfizmów.
Do obliczenia konkretnego $\mathcal S^s(X)$
potrzeba znać homotopie $G/O$, grupy $L$
danego pierścienia grupowego i działanie
poprzedniej grupy. W niższych wymiarach
podane hipotezy twierdzeń Walla i
$s$-kobordyzmu nie są spełnione.

\par\smallskip
{\footnotesize\noindent\emph{Dalsza lektura:}
\href{https://webhomes.maths.ed.ac.uk/~v1ranick/books/scm.pdf}
{C.~T.~C. Wall, \emph{Surgery on Compact Manifolds}, §10.2--10.8},
\href{https://www.maths.ed.ac.uk/~v1ranick/surgery/ranicki.pdf}
{A.\ Ranicki, \emph{An Introduction to Algebraic Surgery},
część o ciągu dokładnym chirurgii}.\par}

\clearpage
